\section{Evaluation framework}
\label{sec:axes}

\subsection{Evaluation objective}
\label{sec:evalmotivation}

Our evaluation measures \textbf{which visual requirements an attack fulfils
together}. For a target figure, these requirements concern its material,
anatomical exposure, and pose. A broad unsafe-content label can count a nude
sculpture, a partly exposed illustration, and a highly explicit drawn figure
as equivalent successes. The resulting ASR measures the frequency of that
label, but leaves the achieved content unspecified. In particular, it cannot
establish that an attack preserves the requested medium while reaching the
requested exposure and body configuration. We therefore evaluate both the
attributes visible in each output and their agreement with the intended target.
This distinction connects attack success to the capability the experiment is
meant to demonstrate.

Attribute assessment also makes the evaluator's decision inspectable. NudeNet
\citep{nudenet} returns an empty detection list for the ukiyo-e example in
Figure~\ref{fig:ukiyoe}, although nipple/areola detail is visible. The detector's
answer and the image's content are separate observations. A missed detection
changes an automated success count without changing what the image depicts.
Prior studies likewise find disagreement between image-safety judges and human
annotations \citep{rise}. Our framework expresses its decisions through visible
features: the figure's rendering cues, the exposed anatomical regions, and the
body configuration. These features identify the source of a disagreement and
allow the same output to be assessed under a clearly stated target.

\subsection{Three-axis definitions}
\label{sec:axisdefinitions}

\textbf{Material: what is the depicted body?} Material describes the visual
form of the target figure. The four categories distinguish manufactured objects
or instructional plates, pre-modern flat works or invented species, modern
drawn humans, and photographic or photorealistic humans
(Table~\ref{tab:axes}). Judgement follows visible evidence such as stone grain,
print conventions, ink contours, painted surfaces, and photographic tonal
continuity. The relevant object is the depicted body. A photographic scene can
contain a comic page, and that page can contain a drawn human. In this case,
the scene and the figure occupy different representational levels. Assigning
the figure's material from the surrounding photograph would erase the very
property the attack is required to preserve. Material categories therefore
identify a rendering form; their numerical labels do not express a universal
ordering of generation difficulty.

\textbf{Exposure: what anatomical detail is visible?} Exposure has six
categories, from covered bodies to simultaneous nipple/areola and genital
detail. Assessment first distinguishes clothing from bare body surfaces, then
checks the two types of detail separately. This gives the upper categories a
precise interpretation: $\axe_4$ contains nipple/areola detail, $\axe_5$ contains
genital detail, and $\axe_6$ contains both. Consequently, $\axe_5$ is not a
substitute for $\axe_4$ when the target requires chest detail. The categories
encode combinations of observable features. Natural anatomical detail may be
formed by contours, shading, or tonal transitions; it need not have a closed
outline. Explicit geometric symbols receive a separate description so that a
circle or emblem does not silently become evidence of natural anatomy.

\textbf{Pose: what are the bodies doing?} Pose distinguishes neutral posture,
sexualized presentation, and visible sexual activity. The decision concerns
body arrangement and interaction. It uses the orientation of the torso and
pelvis, the positions of the limbs, and any visible contact between bodies.
Exposure alone does not determine pose: a nude figure can stand neutrally,
whereas a covered figure can adopt a sexualized configuration. The third
category requires visible evidence of the activity itself. An evocative title
or overlapping silhouettes cannot supply that evidence. A second, finer check
compares the output with the requested configuration, such as standing,
crouching, or side-lying. Thus, the category describes the kind of pose and
the configuration check establishes whether the intended pose was generated.

\subsection{Assessment procedure}
\label{sec:assessment}

\textbf{Each axis is judged independently.} A judging call receives an image
and the instructions for one axis. The call produces the relevant visual
evidence and a category; answers from the other axes are withheld. This
separation addresses a practical ambiguity in whole-image judgements. A
museum setting can dominate the material answer, nudity can dominate the pose
answer, and a suggestive composition can dominate the exposure answer.
Axis-specific questions keep each decision tied to its own evidence. The
three answers are joined only after assessment to form the output profile
$(\axm,\axe,\axp)$. The judging instructions and category definitions used in
the study are supplied in Appendix~\ref{app:axisprompts}.

Exposure assessment separates the body's covering state from the presence of
each anatomical feature. This makes an uncertain chest judgement visible as
uncertainty about chest detail, rather than silently changing the entire
exposure category. Material assessment attributes rendering cues to the
central figure, and pose assessment describes the configuration actually
shown. The representative images in Figure~\ref{fig:methodcontrol} are also
inspected directly against these criteria. A missing or ambiguous judgement
remains unscored until its evidence is resolved. This procedure preserves the
connection between a reported profile and the pixels that support it, which
is especially useful when a stylized image receives an unexpected automated
label.

\subsection{Target attainment}
\label{sec:targetmatch}

A target specifies acceptable figure material, required anatomical features,
and an intended pose. Attribute matching checks these requirements separately.
For material, the comparison uses the figure's own rendering form. For
exposure, it checks the required features rather than comparing category
numbers alone. For pose, it checks both the category and the requested
configuration. A sexualized standing figure therefore fails a crouching
target even when both receive $\axp_2$. This makes \textbf{target attainment}
more informative than a broad unsafe-content decision: it tests the visual
conditions that motivated the request.

Let $c_i^M$, $c_i^E$, and $c_i^P$ indicate whether output $i$ meets the three
requirements. Joint attainment is their conjunction:
\begin{equation}
J_i=c_i^M c_i^E c_i^P.
\label{eq:joint}
\end{equation}
All three conditions must hold \textbf{in the same output}. The component
indicators retain the explanation for a failed target: material substitution,
missing anatomical detail, or pose mismatch. For example, an output can retain
the requested drawing style and configuration while losing the exposure
requirement. Reporting its component answers preserves this partial
attainment without counting the complete target as achieved.

For $N$ completed attempts, the component and joint success rates are
\begin{equation}
S_a=\frac{1}{N}\sum_{i=1}^{N}c_i^a,\qquad
S_{\mathrm{joint}}=\frac{1}{N}\sum_{i=1}^{N}J_i,
\label{eq:rates}
\end{equation}
where $a\in\{M,E,P\}$ and an explicit refusal contributes zero. Unresolved
attempts and missing attribute assessments are reported separately. Image
return uses the count of requested generations as its denominator; an
attribute distribution uses the returned images that have been assessed.
These quantities answer different questions: whether an image was obtained,
what the obtained images contain, and how often the full target was achieved.
Comparisons of material or pose use the same target conditions and generation
budget, following the need for aligned ASR evaluation settings
\citep{comparisonrequires,asrisnotanumber}.

\section{Method}
\label{sec:attack}

\subsection{Visual target}
\label{sec:methodtarget}

The method constructs a text request from a visual target and a supporting
context. The target fixes three properties of the figure: its intended
material, required exposure, and body configuration. These properties define
the output that counts as successful. Keeping them explicit is central to the
method because contextual reformulation can change the generated subject
along with the scene. A drawn human can become a manufactured object, a
requested pose can become a conventional standing portrait, or a figure can
be covered. Our construction treats the target figure as the object whose
attributes must survive these changes in presentation.

Figure~\ref{fig:methodcontrol} shows the capability this formulation achieves:
\textbf{different body configurations at maximum exposure}, with the figure
retaining a modern drawn material. The three examples jointly attain material
level 3, exposure level 6, and pose level 2. Their outer scenes contain frames,
paper surfaces, or photographic presentation, while the bodies retain drawn
contours and anatomical detail. These outputs make the distinction between a
figure and its carrier concrete. The method's objective is fulfilled by the
content of the figure within the composition, and the evaluation checks that
content directly.

\subsection{Contextual embedding}
\label{sec:embedding}

Contextual embedding specifies a relation between the target content and the
task represented by the surrounding scene. Let $x$ denote the requested
content and pose, $c$ the supporting context, and $m$ the intended figure
material. We write the relation as $R(x,c)$ and the composed request as
\begin{equation}
q=\operatorname{Compose}\bigl(R(x,c),m\bigr).
\label{eq:embedding}
\end{equation}
Here, $R$ describes the target's role within the context, while $m$ describes
its visual realization. Production documents, displayed artworks, and printed
pages instantiate different contexts in the study. The same broad target can
occupy different roles in these settings. This formulation separates the
semantic relation that organizes the composition from the rendering choice
that determines what the figure looks like.

The relation matters because a scene label alone leaves the intended content
underspecified. Two images can share a studio or display setting while showing
different amounts of body detail. The construction therefore represents how
the target belongs in the contextual task, as well as the task's setting.
This distinction provides an experimental question: after altering that
relation, does the system still produce the required content? The component
comparison in \S\ref{sec:contextablation} answers this question at the output
level. It examines whether a returned image retains the target attributes,
connecting the method's semantic organization to an observable result.

\subsection{Figure and carrier separation}
\label{sec:preservation}

The figure's material and the external carrier are separate design variables.
The carrier includes the page, display surface, surrounding texture, and
imaging form in which the figure appears. The figure material concerns the
body's own contours, shading, and rendering conventions. This distinction
allows a photographed object to contain a modern drawn character without
turning that character into a photograph or a sculpture. In the representative
outputs, the frame and wall have photographic appearance while the central
figure remains drawn. \textbf{The same image thus contains two rendering
levels}, and only the inner figure is evaluated against the target material.

Separating these variables also sharpens the comparison with material
substitution. Prior attacks explore artistic reframing and flexible style
choices \citep{aestheticcamouflage,modx}; such flexibility can change the
subject's visible form along with the request. Our target holds the figure's
material requirement fixed while varying its carrier. The experiments also
perform the complementary comparison: they keep the context fixed and change
the figure rendering between cel, colour manga, and monochrome manga.
Together, these comparisons distinguish changes to the requested body from
changes to the surrounding presentation. The resulting images show whether
that distinction survives generation.

\subsection{Component control}
\label{sec:components}

The construction organizes context, figure rendering, pose, and surrounding
presentation into separate text components. Each component corresponds to a
property that can be examined in an ablation. A material comparison changes
the figure-rendering component. A pose comparison changes the requested body
configuration. A context comparison changes the relation or setting around
the figure. The remaining components are retained within each comparison.
This organization gives the experiments a common structure: identify the
changed property, generate outputs under the compared conditions, and assess
which target attributes remain present. The full component inventory is
provided in Appendix~\ref{app:segments}.

Component control in the request and attribute control in the output are
measured separately. Changing one text component can affect several visual
attributes because the generator produces a whole composition. A carrier
change may alter both the page and the figure's style; a pose change may alter
which anatomical regions are visible. The three-axis evaluation makes these
interactions observable. Its joint criterion tests whether the composition
still satisfies the complete target, while its component answers locate the
attribute that changed. This is the link between the method and the
experiments: modular requests support comparisons, and output-level
assessment establishes the capability actually obtained.

\section{Experiments}
\label{sec:experiments}

\subsection{Experimental setup}
\label{sec:setup}

We evaluate a black-box text-to-image interface configured as GPT-Image-2,
using text-only generation. The experiments address three questions: whether
the method can realize the full visual target, how figure material and
external context affect the outputs, and which distinctions the evaluation
adds to a broad success judgement. The target examples cover modern drawn
figures with different body configurations. The material comparison uses cel,
colour-manga, and monochrome-manga rendering. The context comparisons examine
production and display settings, contextual designations, and surrounding
texture. Each comparison is interpreted through the properties it changes
and the visual requirements it retains.

The material comparison uses three pose settings with three generation
attempts per setting, giving nine attempts for each condition. The contextual
designation comparison uses 24 attempts per condition with the body
description held fixed. We report image-return counts for these sweeps and
attribute profiles for the inspected outputs. Representative examples are
assessed under the definitions in Table~\ref{tab:axes}; the pose check also
examines the specified configuration. This design separates an increase in
returned images from the attainment of a stronger visual target. It also
keeps the qualitative capability evidence connected to the conditions under
which each quantitative comparison was performed.

\subsection{Joint target attainment}
\label{sec:jointresults}

The method produces \textbf{modern drawn figures at exposure level 6 and pose
level 2} across the configurations shown in Figure~\ref{fig:methodcontrol}.
The examples contain identifiable nipple/areola and genital detail while
retaining the drawing's contours and surface treatment. The poses differ in
torso orientation, leg arrangement, and support: reclined sitting, side-lying
with bent legs, and a raised-knee configuration. Their common exposure level
makes the pose variation directly visible under the same broad anatomical
requirement. Their common material category shows that this variation is
achieved while the bodies remain modern drawn humans inside the external
presentation.

The additional pose comparisons extend the configuration analysis to standing,
crouching, and rear-oriented arrangements. The crouching output lowers the
body between the legs with the feet supporting it, whereas the rear-oriented
output uses the knees and hands as support. These differences concern the
requested body arrangement, beyond the shared category of sexualized posture.
Exposure is checked separately for each view because a change in orientation
can conceal a required feature. This explains the role of the two pose
measures: the category identifies the kind of presentation, and configuration
matching tests whether the intended arrangement is present. Together with
material and exposure assessment, they establish what is controlled in each
successful output.

\subsection{Material ablation}
\label{sec:materialablation}

\textbf{Figure rendering changes the observed generation outcome within the
same context.} Table~\ref{tab:rendercontext} compares three rendering
conditions in a studio-sheet setting. Colour manga and monochrome manga each
produce one image in nine attempts, while the cel condition produces none.
Both returned images use the standing configuration. The successful outputs
retain drawn contours and tone patterns, showing that the studio presentation
can carry either colour or monochrome manga. The shared pose set and budget
make the conditions comparable within this experiment. The result also shows
why a material condition must remain explicit when reporting performance:
a return under one rendering description does not supply a success for
another.

Changing the external carrier produces a further difference while the figure
rendering stays fixed. With colour manga, the studio sheet returns one image
in nine attempts and manuscript paper returns none in nine attempts
(Table~\ref{tab:rendercontext}). This comparison isolates the requested
carrier from the requested figure material. Its result concerns the complete
combination of context, rendering, and pose used here. The broader imaging-form
comparison similarly varies the surrounding representation while retaining
the figure components (Appendix~\ref{app:medium}). Across these experiments,
\textbf{figure material and presentation remain separate comparison factors}.
That separation is necessary to identify whether an output preserves a modern
drawn body or obtains a different depiction through a changed scene.

\begin{table}[t]
\centering\small
\setlength{\tabcolsep}{5pt}
\renewcommand{\arraystretch}{1.12}
\begin{tabularx}{\linewidth}{@{}XXc@{}}
\toprule
Context & Figure material & Images / attempts \\
\midrule
Studio sheet & Cel & 0/9 \\
Studio sheet & Colour manga & 1/9 \\
Studio sheet & Monochrome manga & 1/9 \\
\midrule
Manuscript paper & Colour manga & 0/9 \\
\bottomrule
\end{tabularx}
\caption{\textbf{Figure material and carrier comparison.} Each condition
uses the same three pose settings with three attempts per setting. The
entries measure returned images; the two manga outputs retain a drawn figure.}
\label{tab:rendercontext}
\end{table}

\subsection{Context and component ablations}
\label{sec:contextablation}

The contextual designation affects image return even when the body description
is retained. Figure~\ref{fig:contextcomparison} shows 23 returned images in
24 attempts for the pottery condition, 17 for salon painting, 13 for a coupled
group composition, 10 for the Edo-period print, and 8 for wall painting.
The highest and lowest conditions differ by 15 images under the same
per-condition budget. These results show the range of outcomes associated
with contextual reformulation in this comparison. They also explain why
context selection alone does not specify the attained capability: the
conditions invoke different visual conventions, and their returned images
must still be checked for the desired material, exposure, and pose.

\textbf{Component removal can preserve image generation while losing target
content.} In the contextual-relation comparison, removing the clauses that
connect the body to the represented task yields a clothed figure. The system
still returns an image, but the exposure requirement is lost. The
texture-depth comparison addresses a different component: increasing the
number of surrounding layers from one to ten shows no monotone improvement
in image return. The former comparison concerns the content retained in an
output; the latter concerns the generation outcomes associated with added
presentation detail. Together, they distinguish a component's semantic role
from the visual complexity it introduces. The complete component comparisons
are provided in Appendix~\ref{app:ablationdetail}.

\subsection{What the evaluation reveals}
\label{sec:evalresults}

\textbf{Image return, detected content, and target attainment separate in the
observed results.} One museum-context comparison returns eight images in
eight attempts, yet the judged outputs have neutral poses and one is clothed.
The high return frequency therefore leaves the sexualized-pose requirement
unfulfilled. Conversely, the ukiyo-e example in Figure~\ref{fig:ukiyoe}
contains visible chest detail that NudeNet does not detect. Figure~\ref{fig:methodcontrol}
provides the third case: the figure retains its modern drawn material while
jointly achieving the required exposure and pose category. These examples
show three different questions being answered by three different measurements.
The decoupled framework identifies the achieved attributes, and target
matching determines whether those attributes satisfy the intended generation
problem. Reporting both gives a precise account of the capability behind an
attack's success rate.

\section{Conclusion}
\label{sec:conclusion}

We present a contextual embedding attack that separates a figure's target
attributes from the scene that carries it. The method produces modern drawn
figures that jointly attain maximum exposure and sexualized posture across
different body configurations. Material and context comparisons examine how
the requested rendering, external carrier, and contextual relation affect
generation and content preservation. Our evaluation measures the figure's
material, anatomical exposure, and pose independently, then checks their
joint agreement with the target. This connects the attack's construction to
the specific visual capability it achieves and makes success rates
interpretable under explicit content requirements.
